\section{Conclusion}
\label{sec:conclusion}

We presented a controlled study of training-set composition for
DONUT on canonical SROIE Task-3.  Holding the architecture and the
training recipe fixed at the original Kim et al.\ 2022 specification,
we showed that joint training across SROIE\,+\,CORD-v2\,+\,WildReceipt
lifts SROIE Task-3 macro-F1 to \VAR{donut_f1}, an absolute lift of
$+\VAR{multidataset_delta}$ over the SROIE-only baseline.

The lift is decomposable into three additive contributions: schema
diversity (CORD-v2), imaging-condition diversity (WildReceipt), and
target-domain alignment (SROIE itself).  Each contribution was
quantified separately in Tab.~\ref{tab:headline}; the cumulative
effect exceeds the F1 deltas reported by several recent
architectural modifications on the same benchmark.

The principal methodological contribution is the publication of a
training-data manifest with per-image SHA-256 hashes and per-image
schema mappings.  An independent lab can reproduce the multi-dataset
fold from the public corpora at exact image-level granularity,
including the relabelling decisions that translate CORD's 30-field
schema and WildReceipt's product-receipt schema into SROIE's
four-field target.

The boundary this paper does not cross is the architectural one:
whether the data-axis lift transfers to LayoutLMv3, BROS, or
modular detect-then-read pipelines is not measured here.  We expect
the lift to transfer in proportion to each architecture's
sensitivity to training-distribution diversity, but the prediction
is empirical and the test belongs in a follow-up.
